\documentclass[12pt,a4paper]{article}
\usepackage[utf8]{inputenc}
\usepackage[vietnamese]{babel}
\usepackage{amsmath,amssymb,amsthm}
\usepackage{geometry}
\usepackage{enumitem}
\usepackage[most]{tcolorbox}
\usepackage{hyperref}
\usepackage{fancyhdr}
\usepackage{tikz}
\usepackage{eso-pic}
\usepackage{xcolor}
\geometry{a4paper, margin=2cm, bottom=2.5cm}
\hypersetup{colorlinks=true, linkcolor=blue!70!black}
\setlist[itemize]{topsep=4pt, itemsep=2pt}
\setlist[enumerate]{topsep=4pt, itemsep=2pt}
\AddToShipoutPictureFG{%
  \AtPageCenter{%
    \begin{tikzpicture}[remember picture, overlay]
      \node[rotate=45, scale=1, opacity=0.06]{\fontsize{60}{72}\selectfont\bfseries\sffamily Đặng Tấn Quí};
    \end{tikzpicture}%
  }%
}
\pagestyle{fancy}
\fancyhf{}
\fancyhead[L]{\small\textit{ĐIỀU KHIỂN TỐI ƯU}}
\fancyhead[R]{\small\textit{Đặng Tấn Quí}}
\fancyfoot[C]{\thepage}
\fancyfoot[L]{\footnotesize\textcopyright\ Đặng Tấn Quí}
\fancyfoot[R]{\footnotesize SĐT: 0377\,122\,210}
\renewcommand{\headrulewidth}{0.4pt}
\renewcommand{\footrulewidth}{0.4pt}
\fancypagestyle{plain}{\fancyhf{}\fancyfoot[C]{\thepage}\fancyfoot[L]{\footnotesize\textcopyright\ Đặng Tấn Quí}\fancyfoot[R]{\footnotesize SĐT: 0377\,122\,210}\renewcommand{\headrulewidth}{0pt}\renewcommand{\footrulewidth}{0.4pt}}
\newtcolorbox{dinhnghia}[1][]{enhanced, breakable, colback=blue!3!white, colframe=blue!60!black, fonttitle=\bfseries, title={Định nghĩa}, #1}
\newtcolorbox{dinhly}[1][]{enhanced, breakable, colback=green!3!white, colframe=green!50!black, fonttitle=\bfseries, title={Định lý}, #1}
\newtcolorbox{tinhchat}[1][]{enhanced, breakable, colback=yellow!5!white, colframe=orange!50!black, fonttitle=\bfseries, title={Tính chất}, #1}
\newtcolorbox{phuongphap}[1][]{enhanced, breakable, colback=gray!5!white, colframe=gray!50!black, fonttitle=\bfseries, title={Phương pháp}, #1}
\newtcolorbox{luuy}[1][]{enhanced, breakable, colback=red!3!white, colframe=red!50!black, fonttitle=\bfseries, title={Lưu ý}, #1}
\newtcolorbox{congthuc}[1][]{enhanced, breakable, colback=cyan!3!white, colframe=cyan!50!black, fonttitle=\bfseries, title={Công thức}, #1}
\newtcolorbox{vidu}[1][]{enhanced, breakable, colback=violet!3!white, colframe=violet!50!black, fonttitle=\bfseries, title={Ví dụ}, #1}

\begin{document}
\thispagestyle{empty}
\begin{tikzpicture}[remember picture, overlay]
  \fill[blue!80!black] (current page.south west) rectangle (current page.north east);
  \node[anchor=west, white, font=\fontsize{26}{32}\bfseries\selectfont]
    at ([xshift=3cm, yshift=-7.5cm]current page.north west) {ĐIỀU KHIỂN TỐI ƯU};
  \node[anchor=west, white, font=\fontsize{18}{24}\selectfont]
    at ([xshift=3cm, yshift=-9cm]current page.north west) {Pontryagin, Hamilton--Jacobi--Bellman, LQR, Riccati};
  \node[anchor=west, white, font=\Large\bfseries] at ([xshift=3cm, yshift=-13cm]current page.north west) {Biên soạn:};
  \node[anchor=west, yellow!80!white, font=\fontsize{22}{28}\bfseries\selectfont] at ([xshift=3cm, yshift=-14.5cm]current page.north west) {ĐẶNG TẤN QUÍ};
  \node[anchor=south east, white, font=\Large\bfseries] at ([xshift=-2cm, yshift=3cm]current page.south east) {\the\year};
\end{tikzpicture}
\newpage
\tableofcontents
\newpage
%% ================================================================
%%  CHƯƠNG 1: BÀI TOÁN ĐIỀU KHIỂN TỐI ƯU
%% ================================================================
\section{Bài toán điều khiển tối ưu}

\begin{dinhnghia}[title={Dạng tổng quát}]
  \[
    \min_{u(\cdot)} J[x, u] = \varphi(x(T), T) + \int_{t_0}^{T} L(x(t), u(t), t)\,dt,
  \]
  t.c. $\dot{x}(t) = f(x(t), u(t), t)$, $x(t_0) = x_0$, $u(t) \in U$.

  \begin{itemize}
    \item $x(t) \in \mathbb{R}^n$: trạng thái, $u(t) \in \mathbb{R}^m$: điều khiển.
    \item $\varphi$: chi phí cuối (Mayer), $L$: chi phí chạy (Lagrange), $J$: dạng Bolza.
  \end{itemize}
\end{dinhnghia}

\begin{dinhnghia}[title={Phân loại}]
  \begin{itemize}
    \item \textbf{Thời gian cố định / tự do}: $T$ cho trước hay $T$ cần tối ưu.
    \item \textbf{Trạng thái cuối tự do / ràng buộc}: $x(T) = x_f$ hoặc $\psi(x(T)) = 0$.
    \item \textbf{Ràng buộc điều khiển}: $u \in U$ (tập đóng).
  \end{itemize}
\end{dinhnghia}

%% ================================================================
%%  CHƯƠNG 2: PHÉP TÍNH BIẾN PHÂN
%% ================================================================
\section{Phép tính biến phân}

\begin{dinhnghia}[title={Phiếm hàm và biến phân}]
  $J[y] = \int_a^b F(x, y, y')\,dx$. Biến phân bậc nhất:
  \[
    \delta J = \int_a^b \Bigl(\frac{\partial F}{\partial y}\delta y + \frac{\partial F}{\partial y'}\delta y'\Bigr)dx.
  \]
\end{dinhnghia}

\begin{dinhly}[title={Phương trình Euler--Lagrange}]
  Điều kiện cần: $\frac{\partial F}{\partial y} - \frac{d}{dx}\frac{\partial F}{\partial y'} = 0$.
\end{dinhly}

\begin{tinhchat}[title={Điều kiện biên tự nhiên (Transversality)}]
  Nếu đầu mút tự do: $\frac{\partial F}{\partial y'}\Big|_{x = b} = 0$.
\end{tinhchat}

%% ================================================================
%%  CHƯƠNG 3: NGUYÊN LÝ CỰC ĐẠI PONTRYAGIN
%% ================================================================
\section{Nguyên lý cực đại Pontryagin}

\begin{dinhnghia}[title={Hàm Hamilton}]
  \[
    H(x, u, p, t) = L(x, u, t) + p^T f(x, u, t),
  \]
  $p(t) \in \mathbb{R}^n$: biến đối ngẫu (costate, adjoint).
\end{dinhnghia}

\begin{dinhly}[title={Nguyên lý cực đại Pontryagin}]
  Nếu $(x^*, u^*)$ tối ưu thì tồn tại $p(t)$ sao cho:
  \begin{enumerate}
    \item \textbf{Phương trình trạng thái}: $\dot{x} = \frac{\partial H}{\partial p} = f(x, u, t)$.
    \item \textbf{Phương trình đối ngẫu}: $\dot{p} = -\frac{\partial H}{\partial x}$.
    \item \textbf{Điều kiện cực đại}: $H(x^*, u^*, p, t) \le H(x^*, u, p, t)$ $\forall u \in U$ (min $J$ $\Leftrightarrow$ min $H$ hoặc max $H$ tùy quy ước dấu).
    \item \textbf{Điều kiện biên}: $p(T) = \frac{\partial\varphi}{\partial x}(x(T), T)$ (transversality).
  \end{enumerate}
\end{dinhly}

\begin{luuy}
  Nếu $U = \mathbb{R}^m$ (không ràng buộc): điều kiện cực đại $\Rightarrow$ $\frac{\partial H}{\partial u} = 0$.
\end{luuy}

\begin{vidu}[title={Bang-bang control}]
  $U = [-1, 1]$, $H$ tuyến tính theo $u$: $u^* = -\text{sign}(\frac{\partial H}{\partial u})$. Điều khiển chuyển mạch (switching).
\end{vidu}

%% ================================================================
%%  CHƯƠNG 4: QUY HOẠCH ĐỘNG --- BELLMAN
%% ================================================================
\section{Quy hoạch động --- Phương trình Bellman}

\begin{dinhnghia}[title={Hàm giá trị}]
  \[
    V(x, t) = \min_{u(\cdot)} \Bigl\{\varphi(x(T), T) + \int_t^T L(x, u, s)\,ds\Bigr\}.
  \]
\end{dinhnghia}

\begin{dinhly}[title={Phương trình Hamilton--Jacobi--Bellman (HJB)}]
  \[
    -\frac{\partial V}{\partial t} = \min_{u \in U}\Bigl\{L(x, u, t) + \nabla_x V \cdot f(x, u, t)\Bigr\}, \quad V(x, T) = \varphi(x, T).
  \]
  Nghiệm $V$ cho phép tổng hợp điều khiển phản hồi: $u^*(x, t) = \arg\min_{u}[\cdots]$.
\end{dinhly}

\begin{luuy}
  HJB là PDE phi tuyến bậc nhất; khó giải cho $n$ lớn (\textit{curse of dimensionality}).
\end{luuy}

%% ================================================================
%%  CHƯƠNG 5: ĐIỀU KHIỂN TUYẾN TÍNH --- LQR
%% ================================================================
\section{Điều khiển tuyến tính bậc hai (LQR)}

\begin{dinhnghia}[title={Bài toán LQR}]
  $\dot{x} = Ax + Bu$,
  \[
    J = \frac{1}{2}x(T)^T S_f x(T) + \frac{1}{2}\int_0^T (x^T Qx + u^T Ru)\,dt, \quad Q \succeq 0,\; R \succ 0.
  \]
\end{dinhnghia}

\begin{dinhly}[title={Nghiệm LQR hữu hạn}]
  $u^*(t) = -R^{-1}B^T P(t)x(t)$, với $P(t)$ thỏa phương trình Riccati:
  \[
    -\dot{P} = A^T P + PA - PBR^{-1}B^T P + Q, \quad P(T) = S_f.
  \]
\end{dinhly}

\begin{dinhly}[title={LQR vô hạn (steady-state)}]
  $T \to \infty$, $S_f = 0$. Nếu $(A, B)$ ổn định được thì $P$ nghiệm của phương trình Riccati đại số (ARE):
  \[
    A^T P + PA - PBR^{-1}B^T P + Q = 0.
  \]
  Hệ kín $\dot{x} = (A - BK)x$, $K = R^{-1}B^T P$, ổn định tiệm cận.
\end{dinhly}

%% ================================================================
%%  CHƯƠNG 6: BỔ SUNG
%% ================================================================
\section{Bổ sung}

\begin{tinhchat}[title={Điều khiển được và quan sát được}]
  \begin{itemize}
    \item \textbf{Điều khiển được}: $\text{rank}[B, AB, \ldots, A^{n-1}B] = n$ (Kalman).
    \item \textbf{Quan sát được}: $\text{rank}[C^T, A^TC^T, \ldots, (A^T)^{n-1}C^T] = n$.
  \end{itemize}
\end{tinhchat}

\begin{tinhchat}[title={Ổn định Lyapunov}]
  $V(x) > 0$, $\dot{V}(x) < 0$ dọc quỹ đạo $\Rightarrow$ gốc ổn định tiệm cận.

  Hệ tuyến tính: ổn định $\Leftrightarrow$ tất cả trị riêng của $A$ có phần thực âm.
\end{tinhchat}

\begin{phuongphap}[title={Giải số bài toán điều khiển}]
  \begin{itemize}
    \item \textbf{Shooting}: đoán $p(t_0)$, tích phân hệ trạng thái + đối ngẫu, điều chỉnh.
    \item \textbf{Collocation / Direct}: rời rạc hóa thành bài toán quy hoạch phi tuyến lớn (NLP).
  \end{itemize}
\end{phuongphap}

\end{document}